\subsection{The Dark Sector Crossover}
\label{sec:dark-sector}

The ``Dark Sector''---the sub-lattice of even integers---is geometrically
shielded from the affine frustration of the primes. High-resolution
spectral sweeps ($\Delta N = 2$) across its critical boundary (Exploration~19,
Experiment~B) reveal that its thermalization at $N \approx 150$ is not a
sharp, first-order quantum phase transition, but rather a smooth
Kosterlitz--Thouless-like topological crossover with width
$\Delta N \approx 40$--$60$.

\begin{principle}[Topological Crossover]
The dark sector's smooth thermalization is consistent with a
\emph{topological crossover}: the even integers, connected by a dense
web of divisibility (every pair shares the factor~2), act as a thermal
bath that absorbs the quantum chaos generated by the odd primes through
continuous cross-lattice interactions.

Unlike a true quantum phase transition (which would exhibit a
discontinuous order parameter), the dark sector's GOE fit rises
continuously from $\sim 0.2$ to $\sim 0.8$ as $N$ sweeps through the
critical window $N \in [100, 200]$. The system undergoes a smooth
deconfinement---the even-integer ``bound states'' gradually dissolve
into the chaotic eigenvalue plasma.
\end{principle}

\subsection{The Quantum Decoupling Exponent ($\beta$)}
\label{sec:decoupling}

The spectral decomposition of $d_N^2$ in the eigenbasis of $G_N$
yields the identity
\[
  d_N^2 = 1 - \sum_{k=0}^{N-2} \frac{c_k^2}{\lambda_k},
  \qquad c_k = \langle \mathbf{b}, \mathbf{v}_k \rangle,
\]
where $\lambda_k$ are the eigenvalues and $\mathbf{v}_k$ the eigenvectors
of $G_N$, and $\mathbf{b}$ is the target vector with
$b_j = \int_0^1 \fract{1/(jx)}\,dx$.

The \emph{energy} of the $k$-th mode is $E_k = c_k^2/\lambda_k$,
measuring how much the $k$-th eigenmode contributes to the spectral sum.
For convergence ($d_N^2 \to 0$), the sum $\sum E_k$ must approach~$1$.
The dangerous modes are those with small~$\lambda_k$: if $c_k^2$ does
not decay fast enough relative to~$\lambda_k$, a single mode can cause
$E_k \to \infty$.

\begin{correspondence}[The Quantum Decoupling Exponent]
\label{corr:decoupling}
Define the \emph{decoupling exponent} $\beta$ by the power-law relation
\[
  c_k^2 \sim C \cdot \lambda_k^\beta
\]
fitted over the bottom-$K$ eigenvalues via Lanczos iteration
(50 eigenvalues, Krylov subspace dimension $m = 750$, full
reorthogonalization). The Spectral Observatory
(\texttt{experiments/spectral-observatory/}) measures $\beta$ at
three validated scales:

\begin{center}
\small
\begin{tabular}{@{}rrrrrl@{}}
\toprule
$N$ & $\dim$ & $\beta$ & $\lambda_{\min}$ & $\sum_{k=0}^{49} E_k$ & Source \\
\midrule
10{,}000  & 9{,}999  & 1.611 & $2.54 \times 10^{-7}$ & $8.30 \times 10^{-7}$ & DD MPFR-256 \\
20{,}000  & 19{,}999 & 1.699 & $1.95 \times 10^{-7}$ & $1.25 \times 10^{-6}$ & DD MPFR-256 \\
40{,}000  & 39{,}999 & 1.861 & $1.56 \times 10^{-7}$ & $8.25 \times 10^{-7}$ & DD MPFR-256 \\
\bottomrule
\end{tabular}
\end{center}

The three-point fit yields the scaling law:
\begin{equation}
  \beta(N) \approx -0.062 + 0.180 \cdot \ln(N)
  \label{eq:beta-scaling}
\end{equation}
with residuals $< 0.025$ at all measured points.
$\beta$ crosses $2.0$ near $N \approx 100{,}000$.
\end{correspondence}

\begin{principle}[The Orthogonality Shield]
\label{princ:shield}
The condition $\beta > 1$ means each low-energy mode contributes
vanishing energy: $E_k = c_k^2/\lambda_k \sim \lambda_k^{\beta - 1} \to 0$
as $\lambda_k \to 0$.

The physical mechanism is the \emph{orthogonality shield}: the target
vector $\mathbf{b}$ (which represents the smooth vacuum function~$1$)
has almost zero projection onto the ground-state eigenvectors
$\mathbf{v}_k$, which are localized on composite anchors
(\S\ref{sec:localization}). The projection $c_k = \langle \mathbf{b},
\mathbf{v}_k \rangle$ vanishes because the smooth vacuum is
structurally orthogonal to the spiky, composite-anchored low modes.

As $N$ increases, the Gram matrix becomes increasingly ill-conditioned
(new dangerous modes appear with $\lambda_k \sim N^{-0.35}$), but
the shield \emph{strengthens} logarithmically: $\beta(N) \propto \ln N$.
The arithmetic lattice dynamically segregates its smooth modes from its
chaotic modes, with the segregation growing without bound.
\end{principle}

\begin{remark}[Heisenberg vs.~Schr\"odinger]
The standard proof of $d_N^2 \to 0$ (\texttt{baez\_duarte\_forward},
B\'aez-Duarte 2003) uses the Mellin transform and Parseval identity on
the critical line---the complex-analytic ``Schr\"odinger'' formulation.
The decoupling exponent offers a complementary perspective: $\beta > 1$
is a purely real-spectral condition (Gram eigenvalues and eigenvector
projections), analogous to Heisenberg's matrix mechanics.

$\beta > 1$ is \emph{necessary} for $d_N^2 \to 0$ (it prevents the
bottom modes from diverging) but not \emph{sufficient} alone---the
convergence also requires the bulk spectral mass to accumulate.
The gap is a ``weak completeness'' statement:
\[
  \sum_{\substack{k : \lambda_k \geq \tau}} c_k^2 / \lambda_k
  \;\to\; 1 \quad \text{as } N \to \infty,
\]
which asserts that the spectral weight of $\mathbf{b}$ concentrates
on well-conditioned modes. This is a real-analysis question about
the eigenvalue distribution of $G_N$---potentially tractable via
Weyl-type asymptotics or random matrix universality---rather than
a complex-analytic one.
\end{remark}

\begin{remark}[Precision Requirements]
At $N = 55{,}440$, the MPFR-256 precision (~77 decimal digits) proves
insufficient for spectral extraction: 31 of 50 bottom eigenvalues
appear negative due to roundoff, despite the strict positive
definiteness of~$G_N$. The condition number
$\kappa(G_N) = \lambda_{\max}/\lambda_{\min} > 10^7$ at this scale
exceeds the resolving power of the matrix construction.
MPFR-512 (~154 decimal digits) is required for clean spectral data
at $N > 50{,}000$.

The CG-DD solver for $d_N^2$ survives to higher $N$ because it targets
the quadratic form $\mathbf{b}^T G_N^{-1} \mathbf{b}$ directly
(a scalar, not a spectral decomposition), where error accumulation
is less severe.
\end{remark}

